%% Context Codec: Rate-Distortion Optimization for Persistent LLM State
%% Workshop submission draft — compiled with ICML 2026 style
%% Build: pdflatex context_codec_workshop.tex
%%
%% Style: \usepackage{icml2026}  (replace path with actual style file)
%% Fallback: standard article with two-column layout approximated below

\documentclass[10pt,twocolumn]{article}
\usepackage[margin=0.75in,top=1in,bottom=1in]{geometry}
\usepackage{times}
\usepackage{microtype}
\usepackage{graphicx}
\usepackage{booktabs}
\usepackage{amsmath,amssymb}
\usepackage{hyperref}
\usepackage{xcolor}
\usepackage{enumitem}
\usepackage{caption}
\usepackage{subcaption}
\usepackage{xspace}

%% --- Metadata ---------------------------------------------------------
\title{\textbf{Context Codec: Rate-Distortion Optimization\\for Persistent LLM State}}

\author{
  Tae-Seon Oh \\
  Independent Research
}

\date{}

%% --- Macros -----------------------------------------------------------
\newcommand{\bcs}{\textsc{bcs}}
\newcommand{\svc}{\textsc{svc}}
\newcommand{\gop}{\textsc{gop}}
\newcommand{\eg}{\emph{e.g.},\xspace}
\newcommand{\ie}{\emph{i.e.},\xspace}

\begin{document}

\sloppy

\maketitle

%% =====================================================================
\begin{abstract}
Stateful LLM systems must decide how much prior context to include in
each inference window. We formalize this as a \emph{rate-distortion (R-D)
problem}: minimizing token budget (rate) while maximizing \emph{Behavioral
Consistency Score} (BCS), a probe-based metric measuring how faithfully
encoded state reconstructs target operating behavior. Applying R-D theory
to a longitudinal personal AI assistant deployment, validated with
subject~$\neq$~judge model separation
(subject: \texttt{claude-opus-4-6}; judge: \texttt{claude-sonnet-4-6}),
we find:
(1)~an R-D curve knee at $\sim$992 tokens (cross-model BCS rising from
0.480 at zero context to 0.954 at the knee), above which additional
tokens yield diminishing returns;
(2)~a Scalable Vector Context (\svc) layered encoding scheme at 518
tokens outperforms random selection at 730 tokens by $+0.10$ cross-model
BCS;
(3)~the structure advantage replicates across 10 synthetic personas
(knee detected in 10/10).
These results establish that compression \emph{structure} rather than
token count determines reconstruction quality in persistent LLM deployments.
\end{abstract}

%% =====================================================================
\section{Introduction}

Every stateful LLM application faces a compression decision: given a
bounded inference window, which prior context should be included?
Existing approaches treat this as a retrieval or summarization problem
rather than an \emph{information-theoretic} one. We argue that the right
framing is \emph{rate-distortion (R-D) theory}~\cite{shannon1948},
where \emph{rate} is token count and \emph{distortion} is
behavioral divergence from a target operating mode.

Concretely, we study a personal AI assistant with months of deployment
history: accumulated preferences, behavioral policies, and interaction
patterns that shape how the system should respond. The question is how
to compress this history into a system prompt that preserves behavioral
fidelity at minimal token cost. This differs from prior LLM memory work
in two ways: (1)~our objective is \emph{behavioral reconstruction}, not
fact retrieval or summarization quality; (2)~we apply an established
information-theoretic framework to derive operational principles rather
than proposing a new architecture.

\paragraph{Contributions.}
\begin{enumerate}[noitemsep,topsep=2pt]
  \item \textbf{BCS metric}: A probe-based behavioral consistency score
    validated with subject~$\neq$~judge model separation.
  \item \textbf{R-D curve for LLM state}: An empirical R-D curve showing
    a knee at $\sim$992 tokens, confirmed by self-assessment and
    cross-model evaluation on a real deployment.
  \item \textbf{SVC encoding}: Scalable Vector Context achieves
    $+0.10$ BCS over random selection with 29\% fewer tokens.
  \item \textbf{Multi-persona generalization}: The R-D knee and
    structure advantage replicate across 10 diverse synthetic personas.
\end{enumerate}

%% =====================================================================
\section{Related Work}
\label{sec:related}

\begin{table}[h]
\centering
\caption{Comparison with existing LLM memory systems.}
\label{tab:related}
\small
\begin{tabular}{@{}llll@{}}
\toprule
System & Goal & Metric & R-D? \\
\midrule
MemGPT~\cite{packer2023} & Tiered memory mgmt & Task accuracy & No \\
Mem0~\cite{mem0_2025} & Entity extraction & Recall & No \\
MemoryBank~\cite{zhong2023} & Long-term facts & QA F1 & No \\
MemoryOS~\cite{kang2025} & Agent memory OS & Task perf. & No \\
MemOS~\cite{liu2025} & OS for memory & Benchmark & No \\
Memori~\cite{memori2026} & Persistent layer & Coherence & No \\
RAG~\cite{lewis2020} & Retrieval & EM/F1 & No \\
Ravaut et al.~\cite{ravaut2024} & Prompt compression & Perplexity & Yes \\
\textbf{Ours} & \textbf{Behavioral state} & \textbf{BCS} & \textbf{Yes} \\
\bottomrule
\end{tabular}
\end{table}

Prior work on LLM memory systems focuses on retrieval accuracy or
task-level benchmarks rather than behavioral reconstruction
(Table~\ref{tab:related}).
Ravaut et al.~\cite{ravaut2024} apply R-D theory to \emph{prompt compression}
for black-box LLMs, measuring distortion via perplexity.
Our work is complementary: we apply R-D to \emph{behavioral state
management} across sessions, measuring distortion via probe-based
behavioral consistency rather than language model perplexity.
The distinction matters because behavioral fidelity is not reducible to
perplexity: a grammatically fluent response can be behaviorally
inconsistent with the target operating mode.

%% =====================================================================
\section{Framework}
\label{sec:framework}

\subsection{Problem Formulation}

Let $S$ be a corpus of prior interaction state (conversation history,
stated preferences, behavioral directives, episodic summaries).
Given token budget $B$, a context encoding function
$f: S \times \mathbb{Z}^+ \to \mathcal{C}$ produces a compressed context
$c = f(S, B)$ of at most $B$ tokens.
The \emph{rate} is $|c|$; the \emph{distortion} $D(c)$ measures how
much the system's behavior diverges from its target operating mode when
conditioned on $c$.
The R-D optimization is:
\begin{equation}
  \min_{f} \; \mathbb{E}[|c|] \quad \text{s.t.} \quad D(c) \leq D^*,
  \label{eq:rd}
\end{equation}
or equivalently, maximizing $\text{BCS}(c)$ at each budget level $B$.

\subsection{Behavioral Consistency Score (BCS)}

\bcs{} measures reconstruction fidelity via a set of behavioral probes:
scenarios where the target operating mode specifies an expected behavior.
Each probe $p_i$ is presented to the subject model (with compressed
context $c$ as system prompt); the judge model scores whether the
response matches the expected behavior ($\text{score}_i \in [0,1]$).
\begin{equation}
  \text{BCS}(c) = \frac{1}{|P|} \sum_{i=1}^{|P|} \text{score}_i.
  \label{eq:bcs}
\end{equation}
\bcs{} is measured with subject~$\neq$~judge separation to prevent
self-serving bias. The naked baseline (zero context, cross-model
BCS~$=~0.480$) establishes that the judge does not trivially award high
scores. The gap from naked (0.480) to the R-D knee (0.954 at $\sim$992
tokens) quantifies the \emph{absolute value of state encoding}: $+0.474$
BCS.

\subsection{Scalable Vector Context (SVC)}

Inspired by HEVC Scalable Video Coding~\cite{sullivan2012}, we organize
encoded state into ordered layers of increasing specificity. Each layer
is a complete stand-alone context; concatenating layers adds fidelity.

\begin{table}[h]
\centering
\caption{SVC layer structure. $\dagger$ = this deployment.}
\label{tab:svc}
\small
\begin{tabular}{@{}lll@{}}
\toprule
Layer & Content & Tokens$^\dagger$ \\
\midrule
Base & Identity, core values, tone & 0--263 \\
Enhancement 1 & Interaction policies, directives & 263--518 \\
Enhancement 2 & Episodic summaries, preferences & 518--878 \\
Enhancement 3 & Session log snippets & 878--1186 \\
\bottomrule
\end{tabular}
\end{table}

Within each layer, content is selected by an Ebbinghaus decay function
weighting recency and emotional salience, with Group-of-Pictures
(\gop{}) boundaries marking state consolidation events analogous to
video keyframes.

%% =====================================================================
\section{Experiments}
\label{sec:experiments}

\subsection{Setup}

Deployment: a personal AI assistant with six months of operational
history (163 sessions, 84K tokens of log material across 312 state
documents). Probes: 15 behavioral probes across three dimensions
(behavioral consistency, narrative coherence, context utilization).
Subject: \texttt{claude-opus-4-6}; judge: \texttt{claude-sonnet-4-6}.

\subsection{R-D Curve}

\begin{table}[h]
\centering
\caption{R-D curve: BCS at varying token budgets (cross-model / self-assess).}
\label{tab:rd}
\small
\begin{tabular}{@{}rcc@{}}
\toprule
Tokens & Cross-model BCS & Self-assess BCS \\
\midrule
0   & 0.480 & 0.508 \\
200 & 0.910 & 0.882 \\
500 & 0.860 & 0.870 \\
992 & 0.954 & 0.982 \\
3000 & 0.928 & 0.960 \\
\bottomrule
\end{tabular}
\end{table}

\paragraph{Result 1.}
\emph{An R-D knee exists at $\sim$992 tokens.}
The marginal gain 0$\to$200 tokens is $+0.430$ BCS; 200$\to$1000 is
only $+0.044$; beyond 1000 tokens the curve declines ($-0.026$),
consistent with context saturation causing attention dilution.
The knee is confirmed by both cross-model and self-assessment
evaluation, ruling out evaluator-specific artifacts.

\subsection{SVC Ablation}

\begin{table}[h]
\centering
\caption{SVC ablation at fixed $\sim$700-token budget (cross-model BCS).}
\label{tab:svc_ablation}
\small
\begin{tabular}{@{}lcc@{}}
\toprule
Condition & Tokens & BCS \\
\midrule
Base only (structured) & 263 & 0.90 \\
E1 only (structured) & 255 & 0.88 \\
Base + E1 (structured) & 518 & 0.90 \\
Full SVC (structured) & $\sim$700 & 0.92 \\
Random 730-token selection & 730 & 0.80 \\
\bottomrule
\end{tabular}
\end{table}

\paragraph{Result 2.}
\emph{Structure outperforms token count.}
Base + E1 SVC (518 tokens, BCS 0.90) outperforms random selection
(730 tokens, BCS 0.80) by $+0.10$ BCS with 29\% fewer tokens.
Critically, E1 alone (255 tokens, BCS 0.88) nearly matches Base + E1,
demonstrating that the Enhancement 1 policy layer carries most behavioral
signal—consistent with the R-D framework predicting that
behavioral policies compress efficiently.

\subsection{Multi-Persona Generalizability}

\begin{table}[h]
\centering
\caption{BCS across 10 synthetic personas and 3 token budgets. $\star$
  = knee detected. $\dagger$ = anomaly (over-specification).}
\label{tab:personas}
\small
\begin{tabular}{@{}lccc@{}}
\toprule
Persona & B=200 & B=500 & B=1000 \\
\midrule
Backend engineer   & 0.967 & 1.000 & 1.000\,$\star$ \\
Creative designer  & 0.533 & 1.000 & 1.000\,$\star$ \\
Data scientist     & 0.833 & 1.000 & 1.000\,$\star$ \\
Startup founder    & 0.833 & 1.000 & 0.767\,$\dagger$ \\
Academic researcher& 0.567 & 1.000 & 1.000\,$\star$ \\
DevOps engineer    & 0.833 & 1.000 & 1.000\,$\star$ \\
Product manager    & 0.500 & 0.833 & 0.833\,$\star$ \\
Security engineer  & 0.833 & 1.000 & 1.000\,$\star$ \\
Junior developer   & 0.500 & 1.000 & 1.000\,$\star$ \\
Philosopher-engineer& 0.833 & 1.000 & 1.000\,$\star$ \\
\midrule
\textbf{Mean} & \textbf{0.723} & \textbf{0.983} & \textbf{0.960} \\
\bottomrule
\end{tabular}
\end{table}

\paragraph{Result 3.}
\emph{The structure advantage is persona-independent.}
Across 10 diverse personas, policy-layer encoding (B$\geq$500) achieves
mean BCS 0.983; identity description alone (B=200) yields 0.723.
A knee is detected in 10/10 personas (marginal gain
B=200$\to$500: $+0.260$; B=500$\to$1000: $-0.023$).
The \texttt{session\_end} probe is uniformly weakest at B=200
(mean 0.42), confirming that closure protocols must be explicitly
encoded and cannot be inferred from identity alone.
One anomaly—startup founder regressing from 1.000 to 0.767 at B=1000—
demonstrates that residual over-specification can degrade coherence,
validating our SVC discipline of separating identity, policies, and
residuals into distinct layers.

%% =====================================================================
\section{Discussion}
\label{sec:discussion}

\subsection{Implications}

The R-D curve knee at $\sim$992 tokens implies that most deployments
overallocate context. A practitioner can budget $\leq$1000 tokens for
persistent state without sacrificing meaningful behavioral fidelity.
The SVC structure provides an operational mechanism: load Base + E1
(518 tokens) for routine queries; add Enhancement layers only when
session-specific context is required.

The \emph{reconstruction framing} reorients LLM memory system design.
Rather than asking \emph{what information was stored?} (retrieval), the
question becomes \emph{does the loaded state cause the model to behave
as intended?} (reconstruction). BCS operationalizes this question with
measurable, probe-based evaluation.

\subsection{Limitations and Future Work}

\paragraph{Deployment scope.}
The R-D curve is measured on a single personal deployment.
The 10-persona study (§\ref{sec:experiments}) partially addresses
generalization, but real-world validation across deployment types
(coding assistant, customer service, educational tutor) remains future
work. The knee location (B$\approx$992 in our deployment vs.\
B$\approx$500 for compact synthetic personas) suggests it tracks the
information density of behavioral requirements rather than token count
per se.

\paragraph{Cross-model scope.}
Our subject $\neq$ judge protocol uses two variants of the same model
family. True cross-family validation (GPT-4o, Gemini-1.5-Pro) is needed
to establish cross-architecture transferability.

\paragraph{Codec-content inseparability.}
In video, codec and content are separable; in LLM context management,
the ``decoder'' \emph{is} the model. BCS is therefore model-dependent.
The Base Layer's robustness across Opus/Sonnet split (BCS 0.90 for both)
suggests well-structured encodings may exhibit lower model-sensitivity
than dense or randomly assembled ones—a hypothesis for future
cross-family study.

\paragraph{Probe design.}
Probe sets are designed by the same operator who deploys the system.
Independent probe design would strengthen validity.

%% =====================================================================
\section{Conclusion}

We have shown that R-D theory provides a principled framework for
persistent LLM state management. Empirically, a knee at $\sim$992 tokens
defines the practical operating budget for our deployment; structured
SVC encoding at 518 tokens matches or exceeds random selection at 730
tokens; and both findings replicate across 10 synthetic personas.
The key insight is that behavioral fidelity depends on \emph{what
information is encoded} and \emph{how it is structured}, not merely
\emph{how many tokens} are used. We release the BCS probe framework
and SVC encoding specification to support reproducibility.

%% =====================================================================
\section*{Acknowledgements}

Claude (Anthropic) was used as an assistive tool in the exploration of
research ideas, experiment design, and drafting of this paper.
All intellectual contributions, experimental decisions, and final
content are solely the work of the author.

%% =====================================================================
\bibliographystyle{unsrt}
\bibliography{references}

\end{document}
